Diante disso, operações mais seguras e eficientes demandam menor intervenção do controle terrestre, enquanto a redundância e o uso de algoritmos avançados aumentam a tolerância a falhas. A redução da necessidade de controle remoto, vindos de estações espaciais na Terra, minimiza a latência nas respostas e aprimora a autonomia dos sistemas espaciais, permitindo maior confiabilidade em ambientes complexos e de microgravidade. Além disso, a integração de sistemas de RVD/B com tecnologias autônomas e a fusão de sensores possibilitam a operação em ambientes complexos, como asteroides e a Lua, tudo isso somado permitem o aprimoramento da precisão das manobras, garantindo que acoplamentos mais eficientes sejam realizados e que a segurança dos astronautas sejam estabelecidos em quaisquer condições, mesmo nas mais adversas.
